\documentclass[11pt,a4paper]{article}

\newcommand{\cabecalhodireita}{Lista de Exercícios}
% Preâmbulo compartilhado entre a lista de exercícios e o gabarito.
\usepackage[T1]{fontenc}
\usepackage[utf8]{inputenc}
\usepackage[portuguese]{babel}
\usepackage[a4paper,margin=2.3cm,top=2.0cm,bottom=2.2cm]{geometry}
\usepackage{amsmath,amssymb}
\usepackage{graphicx}
\usepackage{booktabs}
\usepackage{xcolor}
\usepackage{enumitem}
\usepackage{listings}
\usepackage{fancyhdr}
\usepackage[framemethod=default]{mdframed}
\usepackage{titlesec}
\usepackage[hidelinks]{hyperref}

% --- Identidade visual (mesma paleta dos slides: VS Code Light+) ---
\definecolor{ufsverde}{HTML}{1F6F43}
\definecolor{bgcolor}{HTML}{F7F7F7}
\definecolor{linecolor}{HTML}{DDDDDD}
\definecolor{numcolor}{HTML}{999999}
\definecolor{kwcolor}{HTML}{0000FF}
\definecolor{strcolor}{HTML}{A31515}
\definecolor{cmtcolor}{HTML}{008000}

\lstdefinestyle{ufsStyle}{
  backgroundcolor=\color{bgcolor},
  basicstyle=\ttfamily\small\color{black},
  keywordstyle=\color{kwcolor}\bfseries,
  stringstyle=\color{strcolor},
  commentstyle=\color{cmtcolor}\itshape,
  numberstyle=\tiny\color{numcolor},
  numbers=left, numbersep=8pt,
  xleftmargin=18pt, framexleftmargin=18pt,
  breaklines=true, tabsize=4,
  keepspaces=true, showspaces=false,
  showstringspaces=false, showtabs=false,
  frame=single, rulecolor=\color{linecolor},
  framesep=4pt, captionpos=b, upquote=true,
  columns=fullflexible,
  literate={á}{{\'a}}1 {é}{{\'e}}1 {í}{{\'i}}1 {ó}{{\'o}}1 {ú}{{\'u}}1
           {â}{{\^a}}1 {ê}{{\^e}}1 {ô}{{\^o}}1 {ã}{{\~a}}1 {õ}{{\~o}}1
           {ç}{{\c c}}1 {Á}{{\'A}}1 {É}{{\'E}}1 {Í}{{\'I}}1 {Ó}{{\'O}}1
           {Ú}{{\'U}}1 {Â}{{\^A}}1 {Ê}{{\^E}}1 {Ô}{{\^O}}1 {Ã}{{\~A}}1
           {Õ}{{\~O}}1 {Ç}{{\c C}}1,
}
\lstset{style=ufsStyle, language=C}

% Atalho para código no meio do texto
\newcommand{\cd}[1]{\texttt{#1}}

% --- Seções ---
\titleformat{\section}{\normalfont\large\bfseries\color{ufsverde}}{}{0pt}{}
\titlespacing*{\section}{0pt}{14pt}{6pt}

% --- Cabeçalho e rodapé ---
\pagestyle{fancy}
\fancyhf{}
\renewcommand{\headrulewidth}{0.4pt}
\renewcommand{\footrulewidth}{0pt}
\fancyhead[L]{\footnotesize COMP0512 --- Programação C $\cdot$ T02 $\cdot$ 2026.2}
\fancyhead[R]{\footnotesize\textbf{\cabecalhodireita}}
\fancyfoot[C]{\footnotesize\thepage}

% --- Ambiente de exercício ---
\newcounter{exnum}
\newcommand{\exercicio}[1]{%
  \par\refstepcounter{exnum}%
  \vspace{9pt}%
  \noindent{\color{ufsverde}\bfseries Exercício \theexnum.}\hspace{0.4em}\textbf{#1}\par
  \vspace{3pt}%
}

% --- Bloco destacado (instruções, avisos) ---
\newmdenv[
  backgroundcolor=bgcolor,
  linecolor=linecolor,
  linewidth=0.6pt,
  leftmargin=0pt, rightmargin=0pt,
  innerleftmargin=10pt, innerrightmargin=10pt,
  innertopmargin=9pt, innerbottommargin=9pt,
  skipabove=9pt, skipbelow=9pt,
]{destaque}

% --- Cabeçalho de capa ---
\newcommand{\capa}[2]{%
  \noindent
  \begin{minipage}[c]{0.15\textwidth}
    \includegraphics[width=\textwidth]{ufs_principal_positiva.png}
  \end{minipage}%
  \hspace{0.02\textwidth}%
  \begin{minipage}[c]{0.80\textwidth}
    \raggedright\setlength{\parindent}{0pt}
    {\footnotesize Universidade Federal de Sergipe --- Departamento de Computação\par}
    {\footnotesize COMP0512 --- Programação C $\cdot$ Turma T02 $\cdot$ 2026.2\par}
    \vspace{3pt}
    {\Large\bfseries\color{ufsverde} #1\par}
    \vspace{2pt}
    {\small #2\par}
    {\footnotesize Prof.\ André Yoshiaki Kashiwabara \quad$\cdot$\quad \texttt{andreyoshiaki@dcomp.ufs.br}\par}
  \end{minipage}
  \vspace{6pt}
  \hrule height 1pt
  \vspace{10pt}
}


\begin{document}

\capa{Lista de Exercícios --- Fundamentos, Ponteiros e Memória Dinâmica}
     {Conteúdo trabalhado até aqui: fundamentos de C, tipos e TADs, \cd{struct},
      \cd{enum}/\cd{typedef}, ponteiros, vetores e strings, alocação dinâmica de
      memória, vetores e matrizes dinâmicos.}

\section{Parte I --- Fundamentos, tipos e abstração}

\exercicio{Rastreio de laços e divisão inteira.}
Considere o programa abaixo.

\begin{lstlisting}
#include <stdio.h>

int main(void) {
    int n = 5;
    int soma = 0;

    for (int i = 1; i <= n; i++) {
        if (i % 2 == 0) continue;
        soma += i;
    }

    int total = 0;
    for (int i = 0; i < n; i++)
        for (int j = i; j < n; j++)
            total++;

    printf("A: %d\n", soma);
    printf("B: %d\n", total);
    printf("C: %d\n", soma / 2);
    printf("D: %.2f\n", soma / 2.0);
    return 0;
}
\end{lstlisting}

\begin{enumerate}[label=(\alph*),leftmargin=1.8em,itemsep=1pt]
  \item Escreva as quatro linhas impressas.
  \item Quantas vezes o corpo do laço interno (\cd{total++}) executa? Dê a conta,
        não só o número.
  \item As linhas \textbf{C} e \textbf{D} calculam a mesma divisão e imprimem
        valores diferentes. Explique por quê, citando o tipo de cada operando.
\end{enumerate}

\exercicio{Conversões e limites dos tipos.}
Considere o trecho abaixo.

\begin{lstlisting}
char   c = 'A';
int    i = c + 1;
double d = 7 / 2;
double e = 7 / 2.0;
int    f = 3.9;
short  s = 32767;
s = s + 1;

printf("%c %d %.1f %.1f %d %d\n", c, i, d, e, f, s);
\end{lstlisting}

\begin{enumerate}[label=(\alph*),leftmargin=1.8em,itemsep=1pt]
  \item Escreva a linha impressa.
  \item Por que \cd{c + 1} pode ser somado a um inteiro, se \cd{c} é um caractere?
  \item Explique, em uma frase, o valor de \cd{f}: a conversão arredonda ou trunca?
  \item O valor de \cd{s} depende do compilador ou é garantido pela linguagem?
        Justifique e diga como você reescreveria a declaração de \cd{s} para que o
        programa suportasse o valor $32768$.
\end{enumerate}

\exercicio{Tipo, estrutura de dados e TAD: separar o \emph{quê} do \emph{como}.}
Uma \emph{fila de atendimento} de um laboratório guarda, em ordem de chegada, o
nome e a matrícula dos alunos que pediram ajuda. Sobre ela, o resto do programa
precisa poder: criar uma fila vazia, colocar um aluno no fim, retirar o primeiro,
saber quantos estão esperando e saber se está vazia.

\begin{enumerate}[label=(\alph*),leftmargin=1.8em,itemsep=1pt]
  \item Escreva o arquivo de interface \cd{fila.h} \textbf{apenas com} a definição
        dos tipos necessários e os \emph{protótipos} das cinco operações. Não
        escreva nenhum corpo de função.
  \item Descreva, em duas ou três frases, uma implementação possível (o \emph{como}):
        que estrutura de dados você usaria por trás e quais campos ela teria.
  \item Aponte um detalhe da sua implementação que o usuário do TAD \textbf{não}
        precisa conhecer, e explique o que se ganha ao escondê-lo.
\end{enumerate}

\exercicio{Funções sobre um vetor de notas.}
Escreva, em C, as três funções abaixo. Nenhuma delas deve imprimir nada.

\begin{lstlisting}
double media(const double v[], int n);
double maior(const double v[], int n);
int    conta_acima(const double v[], int n, double limite);
\end{lstlisting}

\begin{enumerate}[label=(\alph*),leftmargin=1.8em,itemsep=1pt]
  \item Implemente as três.
  \item Escreva uma \cd{main} que declare \cd{double notas[] = \{7.5, 4.0, 9.25, 6.0, 3.5\};},
        chame as três funções e imprima a média com duas casas decimais, o maior
        valor e quantas notas são maiores ou iguais a $6{,}0$.
  \item Por que o parâmetro é declarado \cd{const double v[]} e não \cd{double v[]}?
        O que o \cd{const} promete a quem chama a função?
  \item O que acontece se alguém chamar \cd{media(notas, 0)}? Altere a sua
        implementação para que esse caso não produza uma operação inválida.
\end{enumerate}

\section{Parte II --- Registros, enumerações e apelidos de tipo}

\exercicio{Mapa de bytes e reordenação de um registro.}
Considere a declaração abaixo, adotando os tamanhos \cd{char}~=~1, \cd{short}~=~2,
\cd{int}~=~4, \cd{float}~=~4, \cd{double}~=~8 e ponteiro~=~8 bytes, com alinhamento
igual ao tamanho do tipo.

\begin{lstlisting}
struct Registro {
    char   sigla;
    double valor;
    int    codigo;
    char   nome[6];
    short  ano;
};
\end{lstlisting}

\begin{enumerate}[label=(\alph*),leftmargin=1.8em,itemsep=1pt]
  \item Some, campo a campo, o tamanho declarado. Quantos bytes dá?
  \item Desenhe o mapa de bytes do registro: uma faixa de 0 até \cd{sizeof(struct Registro)-1},
        marcando onde começa cada campo e onde há \emph{padding}.
  \item Diga o valor de \cd{sizeof(struct Registro)} e de \cd{offsetof(struct Registro, ano)}.
  \item Quantos bytes de \emph{padding} existem no total, e quais são as duas regras
        do compilador que os produziram?
\end{enumerate}

\emph{Ainda sobre a mesma \cd{struct Registro}:}

\begin{enumerate}[label=(\alph*),leftmargin=1.8em,itemsep=1pt,start=5]
  \item Reescreva a declaração reordenando os campos de modo a obter o menor
        \cd{sizeof} possível, sem remover nenhum campo.
  \item Diga o novo \cd{sizeof} e quantos bytes você economizou.
  \item Existe alguma ordenação que chegue a exatamente 21 bytes? Justifique.
  \item Cite uma situação, das discutidas em aula, em que essa economia importa de
        verdade --- e uma em que não vale a pena reordenar.
\end{enumerate}

\exercicio{Registros em funções e vetor de registros.}
Considere o tipo abaixo.

\begin{lstlisting}
typedef struct {
    char   nome[40];
    int    matricula;
    double nota;
} Aluno;
\end{lstlisting}

\begin{enumerate}[label=(\alph*),leftmargin=1.8em,itemsep=1pt]
  \item Escreva \cd{void aplicar\_bonus(Aluno a, double b)} que soma \cd{b} à nota, e
        \cd{void aplicar\_bonus\_ref(Aluno *a, double b)} que faz o mesmo por ponteiro.
        Mostre as duas chamadas a partir de uma variável \cd{Aluno x;} e diga qual
        das duas realmente altera \cd{x}. Explique por quê.
  \item Um registro é copiado inteiro quando passado por valor, mas um vetor não.
        Explique essa diferença em duas frases.
  \item Escreva \cd{int melhor(const Aluno turma[], int n)} que devolve o \emph{índice}
        do aluno de maior nota (em caso de empate, o de menor índice), e
        \cd{-1} se \cd{n <= 0}.
  \item Por que a função devolve o índice em vez de devolver o próprio \cd{Aluno}?
        Cite uma vantagem de cada escolha.
\end{enumerate}

\exercicio{Enumerações: valores, \cd{switch} e o truque do contador.}
Considere a declaração abaixo.

\begin{lstlisting}
typedef enum { NORTE, SUL, LESTE = 10, OESTE, NUM_DIRECOES } Direcao;
\end{lstlisting}

\begin{enumerate}[label=(\alph*),leftmargin=1.8em,itemsep=1pt]
  \item Escreva o valor inteiro de cada uma das cinco constantes.
  \item Um colega escreveu \cd{int contagem[NUM\_DIRECOES] = \{0\};} e pretende usar
        \cd{contagem[d]++} para contar ocorrências de cada direção. O vetor tem
        tamanho suficiente? Há desperdício? Explique o que acontece com esse código.
  \item Reescreva a enumeração de forma que o truque do contador funcione sem
        desperdício, mantendo os quatro nomes de direção.
  \item Escreva uma função \cd{const char *nome\_direcao(Direcao d)} usando \cd{switch},
        com um caso para cada direção. O que o compilador consegue avisar nessa
        função que ele não avisaria se as direções fossem \cd{\#define}?
  \item \cd{Direcao d = 42;} compila? O que isso mostra sobre a garantia que um
        \cd{enum} oferece?
\end{enumerate}

\exercicio{\cd{typedef}: ler declarações e reconhecer armadilhas.}
Responda:

\begin{enumerate}[label=(\alph*),leftmargin=1.8em,itemsep=1pt]
  \item Escreva as \textbf{três} formas de declarar um registro \cd{Ponto} com campos
        \cd{double x, y} usando \cd{struct} e \cd{typedef} (com nome de \emph{tag},
        sem \emph{tag}, e as duas juntas). Diga qual delas permite escrever
        \cd{struct Ponto p;} e qual permite \cd{Ponto p;}.
  \item Dado \cd{typedef double Vetor3[3];}, diga o que é declarado por
        \cd{Vetor3 a;} e o que o parâmetro \cd{void f(Vetor3 v)} realmente recebe.
        Por que essa forma é considerada uma armadilha?
  \item \cd{typedef} cria um tipo novo ou um apelido? Mostre um exemplo curto em que
        essa diferença aparece (isto é, em que o compilador \emph{não} reclama de
        algo que você gostaria que ele reclamasse).
\end{enumerate}

\section{Parte III --- Ponteiros}

\exercicio{Rastreio com \cd{\&}, \cd{*} e aritmética de ponteiros.}
Considere o programa abaixo.

\begin{lstlisting}
#include <stdio.h>

void dobra(int *p, int n) {
    for (int i = 0; i < n; i++)
        *(p + i) = *(p + i) * 2;
}

int main(void) {
    int v[5] = {1, 2, 3, 4, 5};
    int *p = v;
    int *q = &v[3];

    printf("A: %d %d\n", *p, *(p + 2));
    printf("B: %ld\n", (long)(q - p));
    p = p + 1;
    printf("C: %d\n", *p);
    dobra(v, 5);
    printf("D: %d %d %d\n", v[0], *p, *q);
    return 0;
}
\end{lstlisting}

\begin{enumerate}[label=(\alph*),leftmargin=1.8em,itemsep=1pt]
  \item Escreva as quatro linhas impressas.
  \item Suponha que \cd{v} comece no endereço \cd{0x100}. Qual é o valor numérico de
        \cd{p + 2} na linha \textbf{A}? Por que não é \cd{0x102}?
  \item A linha \textbf{B} imprime uma diferença de ponteiros. O que esse número
        significa: bytes ou elementos?
  \item Se a chamada fosse \cd{dobra(p, 4)} em vez de \cd{dobra(v, 5)}, o que a linha
        \textbf{D} imprimiria? Explique.
\end{enumerate}

\exercicio{O bug do \cd{troca} que não troca.}
Um colega escreveu as duas versões abaixo e não entende por que só uma funciona.

\begin{lstlisting}
void troca_v1(int a, int b) {
    int t = a;
    a = b;
    b = t;
}

int main(void) {
    int x = 10, y = 20;
    troca_v1(x, y);
    printf("%d %d\n", x, y);
    return 0;
}
\end{lstlisting}

\begin{enumerate}[label=(\alph*),leftmargin=1.8em,itemsep=1pt]
  \item O que o programa imprime? Explique, falando de \emph{cópias} de argumentos.
  \item Desenhe a memória durante a execução de \cd{troca\_v1}: as variáveis de
        \cd{main} de um lado, os parâmetros da função do outro.
  \item Escreva \cd{troca\_v2} usando ponteiros e mostre a chamada correta a partir
        de \cd{main}. Onde entra o \cd{\&} e onde entra o \cd{*}?
  \item Seu colega tentou \cd{troca\_v2(\&x, \&x)}. O que acontece? A função quebra?
\end{enumerate}

\exercicio{Percorrer sem colchetes.}
Escreva as funções abaixo \textbf{sem usar o operador \cd{[\ ]}} em nenhum lugar ---
apenas ponteiros e aritmética de ponteiros.

\begin{lstlisting}
int  soma_v(const int *v, int n);
void inverte(int *v, int n);
int *primeiro_negativo(int *v, int n);
\end{lstlisting}

\begin{enumerate}[label=(\alph*),leftmargin=1.8em,itemsep=1pt]
  \item Implemente as três. A terceira devolve o \emph{endereço} do primeiro elemento
        negativo, ou \cd{NULL} se não houver nenhum.
  \item Em \cd{inverte}, use dois ponteiros que caminham em sentidos opostos e param
        quando se cruzam. Escreva a condição de parada do laço.
  \item Explique por que \cd{v[i]} e \cd{*(v + i)} são a mesma coisa para o compilador.
  \item Se \cd{primeiro\_negativo} devolvesse o índice em vez do endereço, como o
        chamador distinguiria ``não encontrei''? Compare as duas soluções.
\end{enumerate}

\exercicio{Ponteiro para ponteiro.}
Considere o trecho abaixo.

\begin{lstlisting}
int x = 10, y = 20;
int *p  = &x;
int **pp = &p;

**pp = 30;
printf("A: %d %d\n", x, *p);
*pp = &y;
**pp = 40;
printf("B: %d %d %d\n", x, y, *p);
\end{lstlisting}

\begin{enumerate}[label=(\alph*),leftmargin=1.8em,itemsep=1pt]
  \item Escreva as duas linhas impressas.
  \item Desenhe a cadeia de endereços: \cd{pp} $\to$ \cd{p} $\to$ variável.
        Mostre o desenho \emph{antes} e \emph{depois} da linha \cd{*pp = \&y;}.
  \item Qual é a diferença entre alterar \cd{*pp} e alterar \cd{**pp}?
  \item Cite uma situação prática, vista em aula, em que um parâmetro \cd{int **}
        é necessário.
\end{enumerate}

\section{Parte IV --- Vetores, strings e passagem por referência}

\exercicio{O vetor que encolhe ao entrar na função.}
Considere o programa abaixo.

\begin{lstlisting}
#include <stdio.h>

void mostra(int v[]) {
    printf("dentro: %zu\n", sizeof v);
}

int main(void) {
    int v[10] = {0};
    printf("fora: %zu\n", sizeof v);
    mostra(v);
    return 0;
}
\end{lstlisting}

\begin{enumerate}[label=(\alph*),leftmargin=1.8em,itemsep=1pt]
  \item Escreva as duas linhas impressas e explique a diferença.
  \item Escreva as \textbf{três} declarações equivalentes do parâmetro de \cd{mostra}
        e diga por que as três significam exatamente a mesma coisa para o compilador.
  \item Reescreva \cd{mostra} de modo que ela consiga imprimir quantos elementos o
        vetor tem. O que precisou mudar na assinatura?
  \item Um colega afirmou: ``vetor é ponteiro''. Dê um exemplo curto que mostra que a
        afirmação é falsa.
\end{enumerate}

\exercicio{Strings sem \cd{<string.h>}.}
Implemente à mão, sem usar nenhuma função de \cd{<string.h>}:

\begin{lstlisting}
int  meu_tamanho(const char *s);
void meu_copia(char *destino, const char *origem);
void maiusculas(char *s);
int  meu_compara(const char *a, const char *b);
\end{lstlisting}

\begin{enumerate}[label=(\alph*),leftmargin=1.8em,itemsep=1pt]
  \item Implemente as quatro. \cd{meu\_compara} devolve $0$ se as strings forem iguais,
        um valor negativo se \cd{a} vier antes e positivo se vier depois.
  \item Para \cd{char nome[20] = "Ana";}, diga o valor de \cd{sizeof nome} e o de
        \cd{meu\_tamanho(nome)}. Por que os dois números são diferentes?
  \item Onde está o \cd{'\textbackslash 0'} nesse vetor? Desenhe as 20 posições
        indicando o que se sabe sobre cada uma.
  \item \cd{meu\_copia} tem um risco que \cd{meu\_tamanho} não tem. Qual é, e o que o
        chamador precisa garantir antes de chamá-la?
\end{enumerate}

\exercicio{Devolvendo mais de um resultado.}
Escreva a função

\begin{lstlisting}
void estatisticas(const double v[], int n, double *media, double *menor, double *maior);
\end{lstlisting}

\begin{enumerate}[label=(\alph*),leftmargin=1.8em,itemsep=1pt]
  \item Implemente-a percorrendo o vetor \textbf{uma única vez}.
  \item Escreva a chamada a partir de uma \cd{main} com
        \cd{double v[] = \{3.0, 8.5, 1.25, 6.0\};} e imprima os três resultados.
        Onde aparece o \cd{\&} e por quê?
  \item Por que \cd{v} não precisa de \cd{\&} na chamada, mas \cd{media} precisa?
        Responda usando a regra ``C passa tudo por valor''.
  \item Proteja a função contra \cd{n <= 0} e contra um ponteiro de saída nulo ---
        o chamador que não quiser o maior valor deve poder passar \cd{NULL}.
\end{enumerate}

\exercicio{Dois erros que compilam sem reclamar.}
Analise o programa abaixo.

\begin{lstlisting}
#include <stdio.h>
#include <string.h>

char *saudacao(const char *nome) {
    char buffer[16];
    strcpy(buffer, "Ola, ");
    strcat(buffer, nome);
    return buffer;
}

int main(void) {
    char *s = saudacao("Maria Aparecida da Silva");
    printf("%s\n", s);
    return 0;
}
\end{lstlisting}

\begin{enumerate}[label=(\alph*),leftmargin=1.8em,itemsep=1pt]
  \item Aponte os \textbf{dois} erros distintos e explique o que cada um causa em
        tempo de execução.
  \item Por que o compilador aceita esse código (talvez com um aviso) em vez de
        recusá-lo?
  \item Reescreva \cd{saudacao} de forma correta, \textbf{sem} usar alocação dinâmica.
        Qual foi a mudança na assinatura?
  \item Descreva em uma frase como você reescreveria \cd{saudacao} \emph{com} alocação
        dinâmica, e diga de quem passa a ser a responsabilidade de liberar a memória.
\end{enumerate}

\section{Parte V --- Alocação dinâmica de memória}

\exercicio{Caça às falhas clássicas.}
O programa abaixo compila. Ele contém \textbf{cinco} defeitos ligados ao uso da
memória dinâmica.

\begin{lstlisting}
#include <stdio.h>
#include <stdlib.h>
#include <string.h>

char *duplica(const char *s) {
    char *novo = malloc(strlen(s));
    strcpy(novo, s);
    return novo;
}

int main(void) {
    char *a = duplica("agenda");
    printf("%s\n", a);
    free(a);
    printf("%s\n", a);
    free(a);

    int *v = malloc(10 * sizeof(int));
    v = malloc(20 * sizeof(int));
    free(v);
    return 0;
}
\end{lstlisting}

\begin{enumerate}[label=(\alph*),leftmargin=1.8em,itemsep=1pt]
  \item Liste os cinco defeitos, um por linha, dizendo em que linha cada um está e a
        que categoria pertence (vazamento, uso após liberação, liberação dupla,
        escrita fora do bloco, falta de verificação).
  \item Reescreva o programa corrigido.
  \item O que \cd{free(a)} faz com o \emph{bloco} e o que ele faz com o \emph{ponteiro}
        \cd{a}? Que hábito de escrita evita o defeito das linhas seguintes?
  \item Cite as duas ferramentas vistas em aula para detectar esses defeitos e diga,
        para cada uma, qual dos cinco ela apontaria com mais clareza.
\end{enumerate}

\exercicio{Um vetor que cresce sozinho.}
Implemente uma agenda de inteiros sem limite fixo, com a estrutura e as funções
abaixo.

\begin{lstlisting}
typedef struct {
    int *dados;
    int  tamanho;     /* quantos elementos ja foram inseridos */
    int  capacidade;  /* quantos cabem no bloco atual         */
} Agenda;

int  agenda_cria(Agenda *ag, int capacidade_inicial);
int  agenda_insere(Agenda *ag, int valor);
void agenda_destroi(Agenda *ag);
\end{lstlisting}

\begin{enumerate}[label=(\alph*),leftmargin=1.8em,itemsep=1pt]
  \item Implemente as três funções. \cd{agenda\_insere} deve \textbf{dobrar} a
        capacidade quando o bloco encher, usando \cd{realloc}. As funções que podem
        falhar devolvem $1$ em caso de sucesso e $0$ em caso de falha.
  \item Escreva o \emph{padrão do ponteiro temporário} no \cd{realloc} e explique, em
        duas frases, o que daria errado se você escrevesse
        \cd{ag->dados = realloc(ag->dados, ...)} diretamente e a realocação falhasse.
  \item Depois de um \cd{realloc} bem-sucedido que mudou o bloco de lugar, o que
        acontece com um ponteiro que o chamador tivesse guardado para
        \cd{\&ag->dados[0]}? Como se chama esse problema?
  \item Em \cd{agenda\_cria}, você usaria \cd{malloc} ou \cd{calloc}? Justifique
        dizendo o que muda no conteúdo do bloco --- e por que ``veio zerado no meu
        teste'' não é garantia.
  \item Escreva uma \cd{main} que crie a agenda com capacidade inicial $2$, insira os
        valores de $1$ a $10$, imprima quantas vezes a capacidade mudou e destrua a
        agenda no fim.
\end{enumerate}

\section{Parte VI --- Vetores e matrizes dinâmicos}

\exercicio{Matriz como vetor de ponteiros (\cd{double **}).}
Uma turma tem \cd{nl} alunos e \cd{nc} avaliações, ambos lidos em tempo de execução.

\begin{lstlisting}
double **cria_matriz(int nl, int nc);
void     libera_matriz(double **m, int nl);
\end{lstlisting}

\begin{enumerate}[label=(\alph*),leftmargin=1.8em,itemsep=1pt]
  \item Implemente \cd{cria\_matriz} alocando em \textbf{dois níveis}: primeiro o vetor
        de \cd{nl} ponteiros, depois cada linha com \cd{nc} \cd{double}.
  \item Trate a \textbf{falha parcial}: se a alocação da linha $k$ falhar, desfaça
        exatamente o que já havia sido feito e devolva \cd{NULL}, sem deixar nenhum
        bloco perdido. Escreva esse trecho por inteiro.
  \item Implemente \cd{libera\_matriz} na ordem correta. Um colega escreveu
        \cd{free(m);} antes do laço que libera as linhas. Explique por que isso é um
        erro e que categoria de defeito ele produz.
  \item Desenhe a memória para \cd{nl = 3}, \cd{nc = 4}: o vetor de ponteiros e as três
        linhas. As três linhas ficam necessariamente encostadas na memória?
        Justifique.
\end{enumerate}

\exercicio{A mesma matriz em um bloco só.}
Agora a mesma turma, com \textbf{um único} \cd{malloc}.

\begin{lstlisting}
double *cria_bloco(int nl, int nc);
double  media_aluno(const double *m, int nl, int nc, int i);
double  media_avaliacao(const double *m, int nl, int nc, int j);
\end{lstlisting}

\begin{enumerate}[label=(\alph*),leftmargin=1.8em,itemsep=1pt]
  \item Implemente \cd{cria\_bloco} e escreva a expressão que acessa o elemento da
        linha \cd{i}, coluna \cd{j}.
  \item Implemente \cd{media\_aluno} (média das \cd{nc} notas do aluno \cd{i}) e
        \cd{media\_avaliacao} (média das \cd{nl} notas da avaliação \cd{j}).
  \item Na chamada a \cd{malloc}, por que o material insiste em escrever
        \cd{(size\_t) nl * nc * sizeof(double)} em vez de \cd{nl * nc * sizeof(double)}?
        Descreva a situação em que a diferença aparece.
  \item Um colega trocou os índices e escreveu \cd{m[j * nc + i]}. O programa
        continua compilando e, para uma matriz quadrada, continua rodando sem
        travar. Por que esse é um erro particularmente perigoso?
  \item Compare as duas representações (\cd{double **} e bloco único) em três
        aspectos: número de chamadas a \cd{malloc}, contiguidade na memória e
        facilidade de escrever \cd{m[i][j]}. Diga em que caso você escolheria cada uma.
\end{enumerate}

\end{document}
